% Learning Journey -- Hành Trình Học Thuật Toán
% A visual guide for students on how to learn something completely new
\documentclass[11pt]{article}

% ===== PACKAGES =====
\usepackage[a4paper, margin=2cm]{geometry}
\usepackage[utf8]{inputenc}
\usepackage{fontspec}
\usepackage{amsmath, amssymb}
\usepackage{tikz}
\usetikzlibrary{arrows.meta, positioning, shapes, calc, decorations.pathmorphing, backgrounds, fit, shadows}
\usepackage{tcolorbox}
\tcbuselibrary{skins, breakable}
\usepackage{xcolor}
\usepackage{enumitem}
\usepackage{graphicx}
\usepackage{caption}
\usepackage{hyperref}
\usepackage{multicol}

% ===== COLORS =====
\definecolor{phase1}{HTML}{FF6B6B}   % Red - Overwhelm
\definecolor{phase2}{HTML}{FFA726}   % Orange - Pattern
\definecolor{phase3}{HTML}{66BB6A}   % Green - Build
\definecolor{phase4}{HTML}{42A5F5}   % Blue - Create
\definecolor{phase5}{HTML}{AB47BC}   % Purple - Master
\definecolor{darkbg}{HTML}{1E1E2E}
\definecolor{accent}{HTML}{F9A825}
\definecolor{softgray}{HTML}{F5F5F5}

% ===== CUSTOM ENVIRONMENTS =====
\newtcolorbox{insightbox}[1][]{
  colback=blue!5, colframe=blue!50!black, boxrule=0.8pt, arc=3pt,
  fonttitle=\bfseries, title={#1},
  left=6pt, right=6pt, top=4pt, bottom=4pt
}

\newtcolorbox{storybox}[1][]{
  colback=softgray, colframe=gray!40, boxrule=0.5pt, arc=4pt,
  fonttitle=\bfseries\itshape, title={#1},
  left=8pt, right=8pt, top=6pt, bottom=6pt,
  shadow={1mm}{-1mm}{0mm}{gray!30}
}

\newtcolorbox{mistakebox}[1][]{
  colback=red!5, colframe=red!60!black, boxrule=0.8pt, arc=3pt,
  fonttitle=\bfseries, title={\textcolor{red!60!black}{⚠ #1}},
  left=6pt, right=6pt, top=4pt, bottom=4pt
}

\newtcolorbox{tipbox}[1][]{
  colback=green!5, colframe=green!50!black, boxrule=0.8pt, arc=3pt,
  fonttitle=\bfseries, title={\textcolor{green!50!black}{💡 #1}},
  left=6pt, right=6pt, top=4pt, bottom=4pt
}

% ===== TITLE =====
\title{%
  \vspace{-1cm}
  {\Huge\bfseries Hành Trình Học Một Thứ Mới}\\[0.3cm]
  {\Large\itshape Từ ``Không Biết Gì'' đến ``Hiểu Sâu''}\\[0.4cm]
  {\normalsize Bài học rút ra từ quá trình xây dựng 52 trang Lecture Notes Thuật Toán}
}
\author{Do Minh Phung\\[0.2cm]{\small\itshape Dành cho những ai đang bắt đầu học bất kỳ môn nào mới}}
\date{\today}

\begin{document}
\maketitle

\begin{center}
\begin{tikzpicture}[
  phase/.style={rectangle, rounded corners=6pt, draw=#1, fill=#1!15, 
    text width=2.4cm, minimum height=1cm, align=center, font=\small\bfseries,
    drop shadow={shadow xshift=0.5mm, shadow yshift=-0.5mm, opacity=0.3}},
  arrow/.style={-{Stealth[length=3mm]}, very thick, gray!60},
  label/.style={font=\tiny\itshape, text=gray!70}
]
  \node[phase=phase1] (p1) at (0,0) {Ngợp\\{\footnotesize(Overwhelm)}};
  \node[phase=phase2] (p2) at (3.2,0) {Nhận mẫu\\{\footnotesize(Pattern)}};
  \node[phase=phase3] (p3) at (6.4,0) {Xây dựng\\{\footnotesize(Build)}};
  \node[phase=phase4] (p4) at (9.6,0) {Sáng tạo\\{\footnotesize(Create)}};
  \node[phase=phase5] (p5) at (12.8,0) {Làm chủ\\{\footnotesize(Master)}};
  
  \draw[arrow] (p1) -- (p2);
  \draw[arrow] (p2) -- (p3);
  \draw[arrow] (p3) -- (p4);
  \draw[arrow] (p4) -- (p5);
  
  % Time labels
  \node[label] at (0, -0.9) {Ngày 1--2};
  \node[label] at (3.2, -0.9) {Ngày 3--5};
  \node[label] at (6.4, -0.9) {Tuần 1--2};
  \node[label] at (9.6, -0.9) {Tuần 2--4};
  \node[label] at (12.8, -0.9) {Tháng 2+};
\end{tikzpicture}
\end{center}

\vspace{0.5cm}

% ================================================================
\section{Câu chuyện: Từ ảnh chụp bài giảng → 52 trang Lecture Notes}

\begin{storybox}[Bối cảnh]
Một sinh viên có ảnh chụp bài giảng Thuật Toán, cần tổng hợp thành tài liệu ôn thi. Không biết \LaTeX{}, không rành thuật toán. Trong 2 ngày làm việc, kết quả cuối cùng là \textbf{52 trang Lecture Notes} hoàn chỉnh với:

\begin{multicols}{2}
\begin{itemize}[leftmargin=1.2em, itemsep=2pt]
  \item 20+ pseudocode có annotation
  \item 15+ đồ thị TikZ (cây, graph, heap)
  \item Bảng complexity chi tiết
  \item Animation sắp xếp từng bước
  \item Floyd-Warshall trace đầy đủ
  \item So sánh Greedy vs D\&C vs DP
\end{itemize}
\end{multicols}

\textbf{Quan trọng:} Không phải ``copy-paste'' mà là quá trình \textit{hỏi → hiểu → yêu cầu thêm → phát hiện thiếu → bổ sung}.
\end{storybox}

% ================================================================
\section{5 Giai đoạn Học Một Thứ Mới}

% ----- PHASE 1 -----
\subsection{Giai đoạn 1: Ngợp — ``Cái gì cũng không biết''}

\begin{center}
\begin{tikzpicture}[scale=0.9]
  % Brain with question marks
  \node[circle, fill=phase1!20, draw=phase1, minimum size=2.5cm, line width=1.5pt] (brain) at (0,0) {};
  \node[font=\Huge] at (0, 0.1) {🧠};
  \node[font=\large\bfseries, text=phase1] at (1.6, 0.8) {?};
  \node[font=\Large\bfseries, text=phase1] at (-1.4, 1.0) {?};
  \node[font=\large\bfseries, text=phase1] at (1.2, -0.7) {?};
  \node[font=\large\bfseries, text=phase1] at (-1.6, -0.5) {?};
  \node[font=\normalsize\bfseries, text=phase1] at (0.3, 1.5) {?};
  
  % Labels around
  \node[font=\scriptsize, text=gray] at (3.5, 0.8) {``$\Theta(n \lg n)$ là gì?''};
  \node[font=\scriptsize, text=gray] at (3.8, 0) {``Heap với Stack khác gì?''};
  \node[font=\scriptsize, text=gray] at (4.0, -0.8) {``DP là Dynamic gì?''};
  \node[font=\scriptsize, text=gray] at (-3.5, 0.5) {``Greedy nghe quen''};
  \node[font=\scriptsize, text=gray] at (-3.8, -0.5) {``Floyd ai đó?''};
\end{tikzpicture}
\end{center}

\begin{insightbox}[Đặc điểm của giai đoạn này]
\begin{itemize}[leftmargin=1.2em, itemsep=2pt]
  \item Nhìn đâu cũng thấy \textbf{thuật ngữ lạ} — cảm giác ``mình không thuộc về đây''
  \item Không biết \textbf{bắt đầu từ đâu} — mọi thứ dường như quan trọng như nhau
  \item Dễ bỏ cuộc vì \textbf{khoảng cách} giữa ``hiện tại'' và ``cần đạt'' quá xa
\end{itemize}
\end{insightbox}

\begin{tipbox}[Cách vượt qua]
\textbf{Đừng cố hiểu hết.} Chỉ cần đọc qua 1 lần để biết \textit{có những gì}. Như lần đầu vào thành phố mới -- chỉ cần biết có metro, có bus, có taxi -- chưa cần thuộc tất cả tuyến.

\textbf{Trong dự án này:} Bước đầu tiên chỉ là ``đưa ảnh bài giảng → tổng hợp thành văn bản''. Không yêu cầu TikZ, không yêu cầu đẹp.
\end{tipbox}

% ----- PHASE 2 -----
\subsection{Giai đoạn 2: Nhận mẫu — ``À, hóa ra có pattern!''}

\begin{center}
\begin{tikzpicture}[
  block/.style={rectangle, rounded corners=3pt, draw=phase2, fill=phase2!10,
    minimum width=2cm, minimum height=0.7cm, font=\scriptsize\bfseries, align=center}
]
  % Pattern recognition
  \node[block] (a1) at (0, 1.5) {Bubble Sort};
  \node[block] (a2) at (3, 1.5) {Selection Sort};
  \node[block] (a3) at (6, 1.5) {Insertion Sort};
  
  \node[block, fill=phase2!30, draw=phase2, line width=1.5pt] (pattern) at (3, 0) {Cùng pattern:\\Ý tưởng → Pseudocode\\→ Complexity → Ví dụ};
  
  \draw[-{Stealth}, phase2, thick] (a1) -- (pattern);
  \draw[-{Stealth}, phase2, thick] (a2) -- (pattern);
  \draw[-{Stealth}, phase2, thick] (a3) -- (pattern);
  
  % Lightbulb
  \node[font=\Large] at (8, 0.7) {💡};
  \node[font=\scriptsize\itshape, text=phase2, text width=3cm, align=center] at (8, 0) {``Mỗi thuật toán đều\\có cùng cấu trúc\\trình bày!''};
\end{tikzpicture}
\end{center}

\begin{insightbox}[Đặc điểm của giai đoạn này]
\begin{itemize}[leftmargin=1.2em, itemsep=2pt]
  \item Bắt đầu thấy \textbf{cấu trúc lặp lại}: ``Ờ, sorting nào cũng có Ý tưởng → Pseudocode → Trace''
  \item Bắt đầu \textbf{so sánh}: ``Merge Sort khác Quick Sort ở đâu?''
  \item Bắt đầu \textbf{đặt câu hỏi đúng}: ``Thiếu phần Ý tưởng cho Heap Sort rồi?''
\end{itemize}
\end{insightbox}

\begin{storybox}[Khoảnh khắc ``Aha!'']
Khi nhận ra \textit{``phần Heap Sort thiếu mục Ý tưởng mà các thuật toán khác đều có''} -- đó là dấu hiệu bạn đã nhận ra pattern. Bạn không chỉ đọc, bạn đang \textbf{kiểm tra tính nhất quán} — một kỹ năng của người hiểu sâu.
\end{storybox}

% ----- PHASE 3 -----
\subsection{Giai đoạn 3: Xây dựng — ``Để tôi thử tự làm''}

\begin{center}
\begin{tikzpicture}[
  layer/.style={rectangle, rounded corners=4pt, minimum width=10cm, minimum height=0.9cm, 
    font=\small, align=center, draw=phase3!80},
  >=Stealth
]
  \node[layer, fill=phase3!10] (l1) at (0, 0) {Văn bản thuần túy -- Gõ nội dung};
  \node[layer, fill=phase3!20] (l2) at (0, 1.2) {Thêm bảng, công thức -- $O(n \lg n)$, ma trận};
  \node[layer, fill=phase3!30] (l3) at (0, 2.4) {Thêm pseudocode -- tcolorbox, verbatim};
  \node[layer, fill=phase3!40] (l4) at (0, 3.6) {Thêm đồ thị -- TikZ nodes, edges, trees};
  \node[layer, fill=phase3!55] (l5) at (0, 4.8) {Thêm animation -- ma trận từng bước, trace};
  
  \draw[->, thick, phase3!70] (l1.east) to[bend left=40] node[right, font=\tiny\itshape, text=gray] {+1 kỹ năng} (l2.east);
  \draw[->, thick, phase3!70] (l2.east) to[bend left=40] node[right, font=\tiny\itshape, text=gray] {+1 kỹ năng} (l3.east);
  \draw[->, thick, phase3!70] (l3.east) to[bend left=40] node[right, font=\tiny\itshape, text=gray] {+1 kỹ năng} (l4.east);
  \draw[->, thick, phase3!70] (l4.east) to[bend left=40] node[right, font=\tiny\itshape, text=gray] {+1 kỹ năng} (l5.east);
  
  % Label
  \node[font=\scriptsize\bfseries, text=phase3!80] at (-6, 2.4) {\rotatebox{90}{Độ phức tạp tăng dần →}};
\end{tikzpicture}
\end{center}

\begin{tipbox}[Nguyên tắc vàng: Scaffolding]
Học theo \textbf{từng lớp} (scaffolding), không nhảy cóc:

\begin{enumerate}[leftmargin=1.5em, itemsep=2pt]
  \item \textbf{Làm được} cái đơn giản trước → tự tin
  \item \textbf{Thêm 1 thứ mới} vào cái đã biết → mở rộng
  \item \textbf{Lặp lại} pattern đã học cho trường hợp mới → củng cố
  \item \textbf{Phát hiện giới hạn} → động lực học thêm
\end{enumerate}

\textit{Đừng bao giờ cố học TikZ trước khi viết được 1 đoạn văn bản LaTeX đơn giản.}
\end{tipbox}

% ----- PHASE 4 -----
\subsection{Giai đoạn 4: Sáng tạo — ``Nếu thế này thì sao?''}

\begin{center}
\begin{tikzpicture}[
  idea/.style={cloud, cloud puffs=10, cloud puff arc=120, draw=phase4, fill=phase4!10,
    minimum width=2.5cm, minimum height=1.3cm, font=\scriptsize, align=center, inner sep=2pt},
  >=Stealth
]
  \node[idea] (i1) at (0, 0) {``Thêm ví dụ 3\\cho heap tree?''};
  \node[idea] (i2) at (4.5, 0.5) {``Vẽ đồ thị\\Floyd-Warshall?''};
  \node[idea] (i3) at (9, 0) {``So sánh 3\\phương pháp?''};
  \node[idea] (i4) at (2.5, -1.8) {``Cạnh âm thì\\Dijkstra sao?''};
  \node[idea] (i5) at (7, -1.8) {``Min-Heap vs\\Max-Heap?''};
  
  % Central node
  \node[circle, fill=phase4!30, draw=phase4, minimum size=1.5cm, font=\small\bfseries, line width=1.5pt] (center) at (4.5, -1) {BẠN};
  
  \draw[->, phase4, thick, dashed] (center) -- (i1);
  \draw[->, phase4, thick, dashed] (center) -- (i2);
  \draw[->, phase4, thick, dashed] (center) -- (i3);
  \draw[->, phase4, thick, dashed] (center) -- (i4);
  \draw[->, phase4, thick, dashed] (center) -- (i5);
\end{tikzpicture}
\end{center}

\begin{insightbox}[Dấu hiệu bạn đang ở giai đoạn 4]
Bạn bắt đầu:
\begin{itemize}[leftmargin=1.2em, itemsep=2pt]
  \item \textbf{Đề xuất} thêm nội dung, không chỉ tiếp nhận
  \item \textbf{Phát hiện} lỗ hổng: ``Floyd-Warshall chưa có đồ thị!''
  \item \textbf{Kết nối} kiến thức: ``Dijkstra không chạy với cạnh âm, vậy Floyd-Warshall?''
  \item \textbf{Yêu cầu trực quan}: ``Vẽ cây heap 3 ví dụ song song được không?''
\end{itemize}
\end{insightbox}

% ----- PHASE 5 -----
\subsection{Giai đoạn 5: Làm chủ — ``Tôi có thể dạy lại''}

\begin{center}
\begin{tikzpicture}[
  >=Stealth, scale=0.85, transform shape
]
  % Teaching flow
  \node[rectangle, rounded corners=5pt, fill=phase5!20, draw=phase5, 
    minimum width=3cm, minimum height=1cm, font=\small\bfseries, align=center] (teach) at (0, 0) 
    {Bạn dạy\\cho người khác};
  
  \node[rectangle, rounded corners=5pt, fill=phase5!10, draw=phase5!60, 
    minimum width=2.5cm, minimum height=0.8cm, font=\scriptsize, align=center] (s1) at (-4, 1.5) 
    {Giải thích\\bằng lời mình};
  \node[rectangle, rounded corners=5pt, fill=phase5!10, draw=phase5!60, 
    minimum width=2.5cm, minimum height=0.8cm, font=\scriptsize, align=center] (s2) at (0, 2) 
    {Vẽ hình\\minh họa};
  \node[rectangle, rounded corners=5pt, fill=phase5!10, draw=phase5!60, 
    minimum width=2.5cm, minimum height=0.8cm, font=\scriptsize, align=center] (s3) at (4, 1.5) 
    {Tạo ví dụ\\mới};
  \node[rectangle, rounded corners=5pt, fill=phase5!10, draw=phase5!60, 
    minimum width=2.5cm, minimum height=0.8cm, font=\scriptsize, align=center] (s4) at (-4, -1.5) 
    {Phát hiện\\lỗi sai};
  \node[rectangle, rounded corners=5pt, fill=phase5!10, draw=phase5!60, 
    minimum width=2.5cm, minimum height=0.8cm, font=\scriptsize, align=center] (s5) at (0, -2) 
    {So sánh\\phương pháp};
  \node[rectangle, rounded corners=5pt, fill=phase5!10, draw=phase5!60, 
    minimum width=2.5cm, minimum height=0.8cm, font=\scriptsize, align=center] (s6) at (4, -1.5) 
    {Viết tài liệu\\cho người khác};
  
  \draw[->, phase5, thick] (teach) -- (s1);
  \draw[->, phase5, thick] (teach) -- (s2);
  \draw[->, phase5, thick] (teach) -- (s3);
  \draw[->, phase5, thick] (teach) -- (s4);
  \draw[->, phase5, thick] (teach) -- (s5);
  \draw[->, phase5, thick] (teach) -- (s6);
  
  % Feynman quote
  \node[text width=5cm, align=center, font=\itshape\scriptsize, text=gray] at (8, 0) 
    {``If you can't explain it\\simply, you don't\\understand it well enough.''\\[2pt]--- Richard Feynman};
\end{tikzpicture}
\end{center}

% ================================================================
\section{Biểu đồ: Tiến trình thực tế của dự án}

\begin{center}
\begin{tikzpicture}[
  >=Stealth, scale=0.8, transform shape,
  milestone/.style={circle, fill=#1, draw=#1!80!black, minimum size=5mm, inner sep=0pt},
  mlabel/.style={font=\scriptsize, text width=3cm, align=center},
  time/.style={font=\tiny\itshape, text=gray}
]
  % Timeline
  \draw[very thick, gray!30] (0, 0) -- (16, 0);
  
  % Milestones
  \node[milestone=phase1] (m1) at (0, 0) {};
  \node[milestone=phase1] (m2) at (2, 0) {};
  \node[milestone=phase2] (m3) at (4, 0) {};
  \node[milestone=phase3] (m4) at (6, 0) {};
  \node[milestone=phase3] (m5) at (8, 0) {};
  \node[milestone=phase4] (m6) at (10, 0) {};
  \node[milestone=phase4] (m7) at (12, 0) {};
  \node[milestone=phase5] (m8) at (14, 0) {};
  \node[milestone=phase5] (m9) at (16, 0) {};
  
  % Labels above
  \node[mlabel, above=0.3cm of m1] {Tổng hợp\\ảnh bài giảng};
  \node[mlabel, above=0.3cm of m3] {Thêm bảng\\Complexity};
  \node[mlabel, above=0.3cm of m5] {Vẽ cây\\Heap Sort};
  \node[mlabel, above=0.3cm of m7] {Thêm đồ thị\\Floyd-Warshall};
  \node[mlabel, above=0.3cm of m9] {Bảng so sánh\\3 phương pháp};
  
  % Labels below
  \node[mlabel, below=0.3cm of m2] {Thêm Self-study\\strategies};
  \node[mlabel, below=0.3cm of m4] {Chỉnh sửa\\công thức};
  \node[mlabel, below=0.3cm of m6] {Phát hiện\\thiếu Ý tưởng};
  \node[mlabel, below=0.3cm of m8] {Ví dụ cạnh\\âm YouTube};
  
  % Depth indicator
  \draw[thick, gray!50, dashed] (0, -2.5) -- (16, -2.5);
  \node[font=\tiny, text=gray] at (-1, -2.5) {Nông};
  \node[font=\tiny, text=gray] at (-1, -4.5) {Sâu};
  
  % Depth curve
  \draw[very thick, phase3!80, smooth] plot coordinates {
    (0, -2.7) (2, -3.0) (4, -3.3) (6, -3.6) (8, -3.9) (10, -4.1) (12, -4.3) (14, -4.5) (16, -4.6)
  };
  \node[font=\scriptsize\bfseries, text=phase3!80] at (8, -5) {Độ hiểu sâu tăng dần →};
\end{tikzpicture}
\end{center}

% ================================================================
\section{Bài học rút ra}

\begin{center}
\begin{tikzpicture}[
  lesson/.style={rectangle, rounded corners=5pt, draw=#1, fill=#1!8, 
    text width=6.5cm, minimum height=1.2cm, font=\small, align=left,
    inner sep=6pt},
  num/.style={circle, fill=#1, text=white, minimum size=6mm, font=\small\bfseries, inner sep=0pt}
]
  \node[num=phase1] (n1) at (0, 5) {1};
  \node[lesson=phase1, right=0.3cm of n1] {Bắt đầu bằng \textbf{output thô} -- đừng cầu toàn từ đầu. V1 chỉ cần ``có chữ trên giấy''.};
  
  \node[num=phase2] (n2) at (0, 3.5) {2};
  \node[lesson=phase2, right=0.3cm of n2] {\textbf{Hỏi nhiều} = học nhanh. Mỗi câu hỏi ``tại sao?'' mở ra 1 tầng hiểu biết mới.};
  
  \node[num=phase3] (n3) at (0, 2) {3};
  \node[lesson=phase3, right=0.3cm of n3] {\textbf{Lặp lại pattern} cho nhiều trường hợp. Thêm complexity cho 1 thuật toán → làm cho tất cả.};
  
  \node[num=phase4] (n4) at (0, 0.5) {4};
  \node[lesson=phase4, right=0.3cm of n4] {\textbf{Phát hiện thiếu sót} = dấu hiệu hiểu sâu. ``Heap Sort thiếu Ý tưởng'' → bạn đã maste pattern.};
  
  \node[num=phase5] (n5) at (0, -1) {5};
  \node[lesson=phase5, right=0.3cm of n5] {\textbf{Trực quan hóa} = hiểu cấp cao nhất. Khi bạn vẽ được cây, đồ thị, animation → bạn thực sự hiểu.};
\end{tikzpicture}
\end{center}

% ================================================================
\section{Công thức Học}

\begin{center}
\begin{tikzpicture}
  \node[rectangle, rounded corners=8pt, fill=darkbg, text=white,
    minimum width=13cm, minimum height=3cm, align=center, font=\large,
    inner sep=15pt, drop shadow={shadow xshift=1mm, shadow yshift=-1mm}] {
    \textbf{\textcolor{accent}{Học = Đọc} $\times$ \textcolor{phase1}{Hỏi} $\times$ \textcolor{phase3}{Làm} $\times$ \textcolor{phase5}{Dạy lại}}\\[10pt]
    {\normalsize\itshape
    \textcolor{gray!60}{Đọc 1 lần = 10\% nhớ} \quad|\quad 
    \textcolor{gray!60}{Hỏi ``tại sao'' = 30\%} \quad|\quad
    \textcolor{gray!60}{Tự làm = 70\%} \quad|\quad
    \textcolor{gray!60}{Dạy lại = 95\%}}
  };
\end{tikzpicture}
\end{center}

% ================================================================
\section{Quy trình xác minh đa mô hình AI (Multi-Model Verification)}

Một bài học quan trọng trong dự án này: \textbf{không tin tưởng tuyệt đối vào một AI duy nhất}. Chúng tôi sử dụng nhiều mô hình AI song song, mỗi mô hình có thế mạnh riêng.

\begin{center}
\begin{tikzpicture}[
  >=Stealth, scale=0.85, transform shape,
  model/.style={rectangle, rounded corners=5pt, draw=#1, fill=#1!12, 
    minimum width=3.2cm, minimum height=1.4cm, font=\small, align=center, 
    line width=1.2pt, drop shadow={shadow xshift=0.5mm, shadow yshift=-0.5mm, opacity=0.2}},
  task/.style={rectangle, rounded corners=3pt, draw=gray!50, fill=gray!5,
    minimum width=2.4cm, minimum height=0.6cm, font=\scriptsize, align=center},
  arrow/.style={->, thick, gray!60}
]
  % Models
  \node[model=phase4] (opus) at (0, 0) {\textbf{Claude Opus 4.6}\\{\scriptsize Code · LaTeX · TikZ}\\{\scriptsize System design}};
  \node[model=phase3] (gemini) at (7, 0) {\textbf{Gemini 3.1 Pro}\\{\scriptsize Math logic · Verify}\\{\scriptsize Cross-check}};
  
  % Central workflow
  \node[circle, fill=accent!20, draw=accent, minimum size=1.2cm, font=\scriptsize\bfseries, 
    align=center, line width=1pt] (you) at (3.5, 0) {BẠN\\{\tiny(điều phối)}};
  
  % Arrows
  \draw[arrow, phase4] (you) -- (opus) node[midway, above, font=\tiny] {yêu cầu};
  \draw[arrow, phase3] (you) -- (gemini) node[midway, above, font=\tiny] {verify};
  
  % Tasks for Opus
  \node[task] (t1) at (-2.5, -2) {Viết code LaTeX};
  \node[task] (t2) at (0, -2) {Vẽ TikZ graph};
  \node[task] (t3) at (2.5, -2) {Build pipeline};
  \draw[->, phase4!60, thin] (opus) -- (t1);
  \draw[->, phase4!60, thin] (opus) -- (t2);
  \draw[->, phase4!60, thin] (opus) -- (t3);
  
  % Tasks for Gemini
  \node[task] (t4) at (5, -2) {Verify Floyd trace};
  \node[task] (t5) at (7.5, -2) {Check math logic};
  \node[task] (t6) at (10, -2) {Audit complexity};
  \draw[->, phase3!60, thin] (gemini) -- (t4);
  \draw[->, phase3!60, thin] (gemini) -- (t5);
  \draw[->, phase3!60, thin] (gemini) -- (t6);
  
  % Feedback loop
  \draw[->, red!50, thick, dashed, bend left=25] (gemini.north) to 
    node[above, font=\tiny\itshape, text=red!60] {phát hiện lỗi → sửa lại} (opus.north);
\end{tikzpicture}
\end{center}

\begin{insightbox}[Tại sao cần nhiều AI?]
\begin{itemize}[leftmargin=1.2em, itemsep=3pt]
  \item \textbf{Mỗi model có blind spot.} Model A code giỏi nhưng có thể sai logic toán. Model B kiểm tra logic tốt nhưng không viết được TikZ.
  \item \textbf{Cross-verification} = như có 2 người review code. Nếu cả hai cho cùng kết quả → độ tin cậy cao.
  \item \textbf{Bạn là người điều phối}, không phải người thụ động. Bạn quyết định \textit{hỏi ai}, \textit{hỏi gì}, \textit{khi nào chuyển}.
\end{itemize}
\end{insightbox}

\begin{center}
{\small
\begin{tabular}{|>{\raggedright\arraybackslash}p{3cm}|>{\raggedright\arraybackslash}p{5cm}|>{\raggedright\arraybackslash}p{5cm}|}
\hline
\textbf{Tình huống} & \textbf{Model chính} & \textbf{Model xác minh} \\
\hline\hline
Viết code LaTeX, TikZ & Claude Opus (code mạnh) & Gemini Pro (review output) \\
\hline
Trace thuật toán & Claude Opus (tạo trace) & Gemini Pro (verify từng bước) \\
\hline
Kiểm tra công thức & Gemini Pro (math logic) & Claude Opus (cross-check) \\
\hline
Build pipeline/tool & Claude Opus (system design) & Chạy test thực tế \\
\hline
\end{tabular}
}
\captionof{table}{Phân chia vai trò giữa các mô hình AI}\label{tab:multi-model}
\end{center}

\begin{tipbox}[Quy tắc chuyển model]
\begin{enumerate}[leftmargin=1.5em, itemsep=2pt]
  \item \textbf{Code xong → Verify math.} Sau khi AI viết trace Floyd-Warshall, chuyển sang model khác kiểm tra từng bước $D^{(k)}$.
  \item \textbf{Nghi ngờ → Cross-check.} Nếu kết quả ``trông lạ'', hỏi model thứ hai.
  \item \textbf{Cuối cùng: tự trace bằng tay.} AI verify lẫn nhau chỉ là bước 1. Bước cuối luôn là bạn tự làm.
\end{enumerate}
\end{tipbox}

\vspace{0.5cm}

\begin{mistakebox}[Sai lầm phổ biến]
\begin{enumerate}[leftmargin=1.5em, itemsep=3pt]
  \item \textbf{Đọc xong nghĩ là hiểu.} Hiểu thật = giải thích được cho người khác bằng lời mình.
  \item \textbf{Nhảy thẳng vào cái khó.} Không ai vẽ TikZ graph trước khi viết được ``Hello World'' trong \LaTeX{}.
  \item \textbf{Học một mình.} Hỏi AI, hỏi bạn, hỏi thầy -- mỗi câu hỏi tiết kiệm hàng giờ mò mẫm.
  \item \textbf{Cầu toàn quá sớm.} V1 xấu nhưng \textit{tồn tại} hơn V∞ đẹp nhưng \textit{không bao giờ hoàn thành}.
\end{enumerate}
\end{mistakebox}

\vspace{0.5cm}

\begin{tipbox}[Áp dụng ngay hôm nay]
Chọn \textbf{1 thuật toán} bạn chưa hiểu. Thực hiện đúng 4 bước:

\begin{enumerate}[leftmargin=1.5em, itemsep=3pt]
  \item \textbf{Đọc} pseudocode 1 lần (5 phút)
  \item \textbf{Trace} bằng tay với mảng $[5, 3, 8, 1, 4]$ (10 phút)
  \item \textbf{Vẽ} cây hoặc đồ thị minh họa (10 phút)
  \item \textbf{Giải thích} cho bạn cùng phòng trong 2 phút (2 phút)
\end{enumerate}

\textbf{Tổng: 27 phút.} Nếu bạn làm được 4 bước này, bạn hiểu thuật toán đó hơn 90\% người chỉ đọc slide.
\end{tipbox}

\vspace{1cm}
\begin{center}
\rule{0.5\linewidth}{0.5pt}\\[0.3cm]
{\itshape ``Hành trình vạn dặm bắt đầu từ một bước chân.''}\\[0.2cm]
{\small --- Lão Tử (老子)}
\end{center}

\end{document}
